\section{Introduction}

Multi-agent systems designed for long-horizon task execution face a fundamental reliability hazard we term the \textit{autonomous drift problem}. As agents recursively spawn sub-agents to handle partitioned workloads, the resulting hierarchy undergoes unintended structural evolution: agents diverge from their originating directives, inherit conflicting or stale context, or propagate goals along chains that progressively distort the original intent. The longer the system runs and the deeper the spawning hierarchy grows, the further any given agent can stray from what the system operator actually intended. In production deployments this manifests as agents that continue executing tasks long after their parent objectives have been superseded, consume compute and budget on goals that no longer map to any live business objective, or silently abandon critical paths because a distant ancestor was modified or deleted without notifying its descendants.

Existing approaches to containing this failure mode fall into two categories, both of which we argue are structurally insufficient for high-stakes, long-running deployments. The first category is temporal bounding via time-to-live (\texttt{TTL}) mechanisms. Agents or their tasks are assigned an expiration timestamp, and the system terminates any agent whose \texttt{TTL} has elapsed. This approach is intuitive and cheap to implement, but it is fundamentally a blunt instrument: it kills agents based on \textit{age}, not on whether their current activity remains aligned with live system objectives. An agent whose work is still relevant but whose \texttt{TTL} has expired is terminated indiscriminately, while an agent that has already drifted far from its mandate but was spawned more recently continues executing unimpeded. \texttt{TTL} bounds also require operators to correctly anticipate execution duration a priori, which is impossible for adaptive, exploratory, or conditionally branching workflows. The second category is depth-limiting: spawning hierarchies are capped at a fixed maximum depth, beyond which no further sub-agents may be created. Depth limits introduce a different failure mode: they do not prevent drift within the capped region, and they create an arbitrary structural wall that either throttles parallelism prematurely or forces the system to route complex tasks through a fixed number of shallow agents, degrading the system's ability to handle tasks whose natural decomposition requires variable depth.

A closer examination reveals that both failure modes share a common root cause. In conventional multi-agent architectures, when an agent spawns a sub-agent, the sub-agent receives a reference to its parent---or to a shared task context---but this reference is \textit{mutable}. The parent may be reclassified, its priority changed, its context updated, or the parent agent itself may be terminated or replaced. Because the spawning chain is mutable, any agent that retains a reference to an ancestor is holding a handle to a reality that can change beneath it without notice. Drift occurs not because agents behave maliciously, but because the structural invariants of the hierarchy itself are fragile: there is no enforcement mechanism that guarantees an agent's ancestry remains coherent and traceable from inception to termination.

We propose a structurally different approach grounded in a single, deceptively simple invariant: \textbf{immutability of the parent pointer chain}. In the \texttt{BX3} framework, every agent maintains an immutable, append-only chain of parent pointers that records its complete spawning ancestry at creation time. This chain is sealed at the moment of agent instantiation and cannot be altered throughout the agent's lifetime, nor can it be altered in any ancestor after the fact. The critical consequence is that an agent's spawning lineage is a fixed, verifiable record. There is no mutable shared state to diverge from, no parent reference that can be silently rewritten, and no dangling or updated ancestry that an active agent might be stale with respect to. Drift, in this model, becomes structurally impossible: an agent cannot meaningfully diverge from its ancestry because its ancestry is definitionally immutable and its execution is anchored to that fixed chain.

We sharpen this intuition with a formal definition. Let an agent $a_i$ be instantiated at time $t_i$ with an immutable parent pointer $p(a_i)$ that references the agent that spawned it, or \texttt{nil} if it is a root agent. Recursively applying the parent pointer operator yields the \textit{spawning ancestry chain} $\mathcal{A}(a_i) = (a_i, p(a_i), p(p(a_i)), \dots)$ terminated at the root. The chain is sealed at $t_i$; no operation in the system may alter any element of $\mathcal{A}(a_i)$ after $t_i$. We define \textit{drift} as the condition in which the effective objective of agent $a_i$ at time $t > t_i$ is not contained within the transitive closure of the objectives of the agents in $\mathcal{A}(a_i)$ at their respective instantiation times. In conventional mutable-chain systems, drift is possible because $p(a_i)$ is dynamic and the objectives of ancestors may change post-instantiation. In the \texttt{BX3} immutability model, drift is ruled out by construction: the ancestry chain is frozen, and any cross-agent coordination must proceed through explicit, auditable message passing rather than implicit shared state mutation.

This paper is organized as follows. In Section~\ref{sec:bg} we situate \texttt{BX3} within the landscape of agentic architectures and multi-agent coordination paradigms. In Section~\ref{sec:model} we present the formal agent model, including the immutable parent pointer chain, drift definition, and the spawn operation semantics. In Section~\ref{sec:coordination} we describe the coordination protocols that become possible when ancestry is immutable: traceable error propagation, deterministic teardown, and non-equivocating inter-agent contracts. In Section~\ref{sec:implementations} we detail two reference implementations---a local inference variant using Ollama and a cloud-native variant using the Mistral API---and report on empirical evaluation measuring drift resistance, latency, and resource efficiency. In Section~\ref{sec:related} we examine related work in agent architectures, spawning models, and multi-agent reliability. We conclude in Section~\ref{sec:conc} with a discussion of limitations and a roadmap for extending the \texttt{BX3} model to support conditional ancestry branches and adversarial robustness tests.
